\id{ҒТАМР 65.63.29}{}

\begin{header}
\swa{Ш.Ж.~Жасқайрат, Н.С.~Машанова, Д.С.~Жаксыгалиева, М.Е.~Бекболатова}{ЛАКТОЗАСЫ ТӨМЕН ФУНКЦИОНАЛДЫҚ СҮТТІ СУСЫННЫҢ ТЕХНОЛОГИЯСЫН ӘЗІРЛЕУ}

\tsp{1}Ш.Ж.~Жасқайрат\envelope,
\tsp{1}Н.С.~Машанова,
\tsp{2}Д.С.~Жаксыгалиева,
\tsp{1}М.Е.~Бекболатова
\end{header}

\begin{affil}
\tsp{1}С.Сейфуллин атындағы Қазақ агротехникалық зерттеу университеті, Астана, Қазақстан,

\tsp{2}Жәңгір хан атындағы Батыс Қазақстан аграрлық-техникалық университеті, Орал, Қазақстан

\corrauthor{Корреспондент-автор: shynarai\_92@mail.ru}
\end{affil}

Бұл мақалада адамның дұрыс тамақтану қажеттілігіне сәйкес лактозасы
төмен функционалдық сүт өнімін өндіру технологиясын талдау нәтижелері
көрсетілген. Талдау нәтижелері бойынша дайын сүт өнімінің
органолептикалық және физика-химиялық көрсеткіштері сиыр сүтінің
көрсеткіштерімен салыстыра отырып анықталды, яғни сыртқы түрі біркелкі,
консистенциясы біртекті, иіс пен дәм бойынша сүтке тән дәм, ащы емес,
біркекі ақ түске ие. Дайын өнімнің физика-химиялық көрсеткіштері:
тығыздығы - 28,77 \(кг\)/м\tsp{3}, қышқылдығы - 19 °Т, ақуыз
-- 3,09 г, лактоза -- 0,1 г. Сиыр сүтіне лактаза ферменті мен әртүрлі
мөлшердегі теңіз балдыры қосылып, төрт үлгі дайындалды. Зерттеу
нәтижесінде 0,02 \% теңіз балдыры қосылған сүт үлгісі органолептикалық
және физика-химиялық көрсеткіштері бойынша оңтайлы деп танылды.

{\bfseries Түйін сөздер}: сүт, лактоза, фермент, теңіз балдыры,
физика-химиялық көрсеткіштер, органолептикалық көрсеткіштер.

\begin{header}
РАЗРАБОТКА ТЕХНОЛОГИИ ФУНКЦИОНАЛЬНОГО МОЛОЧНОГО НАПИТКА С НИЗКИМ СОДЕРЖАНИЕМ ЛАКТОЗЫ

\tsp{1}Ш.Ж.~Жасқайрат\envelope,
\tsp{1}Н.С.~Машанова,
\tsp{2}Д.С.~Жаксыгалиева,
\tsp{1}М.Е.~Бекболатова
\end{header}

\begin{affil}
\tsp{1}Казахский агротехнический исследовательский университет им. С.Сейфуллина, Астана, Казахстан,

\tsp{2}Западно-Казахстанский аграрно-технический университет, Уральск, Казахстан,

e-mail: shynarai\_92@mail.ru
\end{affil}

В данной статье отражены результаты анализа технологии производства
малофункциональной молочной продукции с низкой лактозой в соответствии с
потребностями человека в здоровом питании. По результатам анализа
органолептические и физико-химические показатели готовой молочной
продукции определены в сравнении с показателями коровьего молока, то
есть имеют равномерный внешний вид, однородную консистенцию, характерный
для молока вкус по запаху и вкусу, не горький, а несколько белый.
Физико-химические показатели готовой продукции: плотность - 28,77
кг/м\tsp{3}, кислотность - 19 ° Т, белок - 3,09 г, лактоза -
0,1 г. К коровьему молоку добавлены фермент лактазы и различные дозы
морских водорослей, изготовлено четыре образца. В результате
исследования образец молока с добавлением 0,02 \% морских водорослей был
признан оптимальным по органолептическим и физико-химическим
показателям.

{\bfseries Ключевые слова:} молоко, лактоза, фермент, морские водоросли,
физико-химические показатели, органолептические показатели.

\begin{header}
DEVELOPMENT OF A FUNCTIONAL LOW-LACTOSE DAIRY BEVERAGE TECHNOLOGY

\tsp{1}Sh.Zh.~Zhaskairat\envelope,
\tsp{1}N.S.~Mashanova,
\tsp{2}D.S.~Zhaksygalieva,
\tsp{1}M.E.~Bekbolatova
\end{header}

\begin{affil}
\tsp{1}Saken Seifullin Kazakh Agrotechnical Research University, Astana, Kazakhstan,

\tsp{2}Khan Zhangir West-Kazakhstan agrarian technical University, Uralsk, Kazakhstan,

e-mail: shynarai\_92@mail.ru
\end{affil}

This article presents the results of the analysis of the production
technology of low-functional dairy products with low lactose in
accordance with the human need for healthy nutrition. Based on the
results of the analysis, the organoleptic and physicochemical indicators
of the finished dairy products were determined in comparison with the
indicators of cow' s milk, that is, they have a uniform
appearance, homogeneous consistency, a characteristic milk taste in
smell and taste, not bitter, but somewhat white. Physicochemical
indicators of the finished product: density -- 28,77 kg /
m\tsp{3}, acidity - 19 ° T, protein -- 3,09 g, lactose --
0,1 g. Lactase enzyme and various doses of seaweed were added to
cow' s milk, four samples were made. As a result of the
study, the milk sample with the addition of 0,02 \% of seaweed was
recognized as optimal in terms of organoleptic and physicochemical
indicators.

{\bfseries Key words}: milk, lactose, enzyme, seaweed, physicochemical
parameters, organoleptic parameters.

\begin{multicols}{2}
{\bfseries Кіріспе.} Адамның дұрыс тамақтануы үшін құрамы теңдестірілген,
өсімдік ингредиенттерімен байытылған функционалды сүт өнімдерін өндіру
технологиясын жасау өзекті мәселе болып табылады. Функционалды тамақ
өнімдері, соның ішінде лактозасы аз және лактозасыз тамақ өнімдері
құрамында физиологиялық функционалды тағамдық ингредиенттердің болуына
байланысты адамның денсаулығын сақтауға, жақсартуға және аурулардың
алдын алуға көмектесетін пайдалы тағам өнімі {[}1{]}.

Адамдардың едәуір бөлігі ферментопатияның көбеюіне байланысты сүт
өнімдерін тұтына алмайды. Төмен лактозалы сүт өнімдерінің құрамында
лактоза мөлшері аз болуымен қатар, тағамдық сипаттамалары бойынша
кәдімгі лактозасы бар сүт өнімдеріне ұқсас болуы тиіс. Сондықтан осындай
тағам өнімдерін өндірушілер тұтынушылардың қажеттіліктері мен қабылдау
ерекшеліктерін мұқият зерделеуі қажет.1-суретте қазіргі әлемдік нарықта
төмен лактозалы өнімдердің көптеген түрлері ұсынылған: ішуге арналған
сүт, сүтті өнімдер, ірімшік, балмұздақ, сүт консервілері және т.б.
өнімдер көрсетілген.
\end{multicols}

\fig{p/image47}[1-сурет. Төмен лактозалы өнімдердің
өнім түрі бойынша ассортимент құрылымы, \%]

\begin{multicols}{2}
Соңғы уақытта төмен лактозалы және лактозасыз сүт өнімдеріне сұраныс
әлемдік деңгейде қарқынды өсуде. Мысалы,~Dairy Global (2023)~деректері
бойынша төмен лактозалы өнімдер нарығы жыл сайын шамамен~7-9 \%өсіп
отыр, ал~Еуропа мен Солтүстік Америка~елдерінде мұндай өнімдердің үлесі
жалпы сүт өнімдерінің~25 \%-ына~жетті.

Қазіргі уақытта сүт шикізатындағы лактоза мөлшерін төмендету үшін
әртүрлі әдістер қолданылуда. Олардың ішінде: лактаза ферментін пайдалану
арқылы ферментативті гидролиз; баромембраналық сүзу; хроматографиялық
тазарту; сүттің түрлі микро және макрокомпоненттерін сарысу ақуызы,
казеин немесе казецитпен біріктіру және т.б. әдістер {[}2- 4{]}.

Шетелдік ғалымдардың зерттеулерінде төмен лактозалы өнімдерді
өндіруде~β-галактозидаза ферментін пайдалану, сондай-ақ~мембраналық
сүзу~және~ферментативтік гидролиз~әдістері тиімді деп танылған. Мысалы,
Knudsen L.J. et al. (2024)~зерттеуінде гидролиз әдісінің түрі мен
ферменттің табиғаты лактозаны ыдырату тиімділігіне және өнімнің дәмдік
қасиеттеріне елеулі әсер ететіні дәлелденген {[}5{]}.

Titov E.I. және әріптестері (2020)~лактаза типіне байланысты гидролиз
тиімділігін зерттеп, β-галактозидаза қолдану арқылы лактозаны 98 \%-ға
дейін азайтуға болатынын көрсетті {[}6{]}.

O' Sullivan A.M. et al. (2014)~теңіз балдырларының сүт
өнімдерінің функционалдық қасиеттерін күшейтетін табиғи ингредиент
екенін дәлелдеді {[}7{]}.

Orzuna-Orzuna J.F. және т.б. (2024)~зерттеуінде теңіз балдырларын сүт
сиырларының рационына енгізу сүт өнімдерінің сапалық құрамын арттыратыны
және экологиялық әсерін төмендететіні көрсетілген {[}8{]}.

Сонымен қатар,~Mummah S. et al. (2014)~зерттеуінде шикі және төмен
лактозалы сүтті тұтынудың лактозаға төзбеушілік белгілерін айтарлықтай
төмендететіні анықталған {[}9{]}.

Осы халықаралық зерттеулер төмен лактозалы өнімдердің физиологиялық және
технологиялық артықшылықтарын дәлелдейді және бұл бағыттағы
ғылыми-зерттеу жұмыстарын одан әрі дамытудың өзектілігін арттырады.

Зерттеудің мақсаты - лактозасы төмен функционалды сүт сусынын өндіру
технологиясын әзірлеу және теңіз балдырлары қосылған өнімнің
органолептикалық және физика-химиялық қасиеттерін бағалау.

Зерттеудің гипотезасы: теңіз балдырларын аз мөлшерде (0,02 \%) қосу
өнімнің құрылымдық-тұтқырлық қасиеттерін және тағамдық құндылығын
жақсартады, ал лактаза ферментін пайдалану сүттегі лактоза мөлшерін
едәуір төмендетеді.

{\bfseries Материалдар мен әдістер.} Зерттеу нысаны ретінде С.Сейфуллин
атындағы Қазақ агротехникалық университетінің «Сүтті қайта өңдеу
эксперименттік-өндірістік цехында» өндірілген сиыр сүтінен сүт сарысуы
алынды. Зерттеу нысандары:

- стерильденген сүт, майлылығы 2,5\%;

- пастерленген сүт, майлылығы 3,9\%;

- балғын сүт;

- лактаза ферменті

Жұмыс барысында қолданылатын барлық компоненттер қолданыстағы
нормативтік-техникалық құжаттаманың талаптарына сәйкес.

Бұл жұмыста шикізаттар мен дайын өнімнің қасиеттерін зерттеу үшін
органолептикалық, физико-химиялық зерттеу әдістері қолданылды.

1. Сүттің органолептикалық қасиеттерін анықтау МЕМСТ 28283-2015
стандартына сайкес анықталды.

Сүттің түсін анықтау. Өлшеу цилиндрінде әр типтегі 50 мл сүт өлшеніп,
химиялық шыныға құйылды.Ақ фондық экранды қолдана отырып, әр үлгінің
түсі анықталды.

Сүттің иісін анықтау. Әр үлгі 20 мл су моншасында 35
\tsp{0}С температураға дейін қыздырылды.

Сүттің дәмін анықтау.1, 2, 3 және 4 үлгідегі сүттің дәміне бағалау
жүргізілді. Шикі сүттің дәмін бағалау (3) санитарлық-гигиеналық
ережелерге сәйкес жүргізілді.

2. Сүттің физика-химиялық қасиеттерін анықтау.

Сүт және сүт өнімдерінің тығыздығын анықтау. Сүттің тығыздығы МЕМСТ Р
54758-2011 стандартына сай ареометрде анықталды. Сүт бөлме
температурасына дейін қыздырылып, біркелкі араластырылды. Цилиндрге 250
мл сүт құйылып, сүт температурасы өлшенеді. Ареометр сүтке батырылып,
еркін жүзу күйінде қалдырылады да ареометр шкаласы бойынша тығыздық
көрсеткіші анықталады.

3. Сүт және сүт өнімдерінің қышқылдығын анықтау. Сүттің қышқылдығы МЕМСТ
Р 54669-2011 стандартына сай анықталды. Бюретка нарий гидроксидінің
ерітіндісімен (0,1н) толтырылып, оның деңгейін нөлдік бөліністе орнатты.
Сыйымдылығы 100 мл конустық колбаға 10 мл сүт өлшеніп, үстіне 20 мл
дистилденген су мен үш тамшы фенолфталеин тамызылды. Қоспа мұқият
араластырылып, колбаны сол қолмен ұстап, оның құрамын жеңіл айналмалы
қозғалыспен үздіксіз араластыра отырып, 1 минут ішінде жоғалып кетпейтін
сәл қызғылт бояу пайда болғанға дейін 0,1 н натрий гидроксиді
ерітіндісімен бір тамшыдан тамызылды. Тернер градусындағы сүттің
қышқылдығын анықтау үшін 10 мл сүтті титрлеуге жұмсалған сілтінің (мл)
көлемі 10-ға көбейтілді, яғни 100 мл өнімге қайта есептелді.

4. Сүт және сүт өнімдерінің құрамындағы ақуызды анықтау МЕМСТ 34536-2019
стандартына сай жүргізілді. Колбаларға 5 мл сүт құйылып, әрқайсысына 5
тамшы сірке қышқылы ерітіндісі тамызылды. Әрі қарай оны 40-45 °С
температураға дейін су моншасында қыздырып, центрифугада (8000-10000
айн./мин) 10 минут бойы айналдырды. Нәтижесінде сүттің тұнбасы байқалды,
содан кейін ерітінді қағаз сүзгісі арқылы сүзілді. Сүзгіде сүзбе тұнбасы
қалды, ол пробиркаға салынып, 3 тамшы концентрацияланған азот қышқылы
тамызылды. Сәйкесінше ақуыздардың тұнбада болуы ашық сары түстің пайда
болуымен анықталды.

5. Сүт және сүт өнімдерінің құрамындағы лактозаны анықтау МЕМСТ
34304-2017 стандартына сай анықталды. Алдыңғы тәжірибеде алынған
фильтрат фенолфталеин индикаторы бойынша натрий гидроксиді ерітіндісімен
бейтараптандырылды. Фильтрат екі жартыға бөлінді. Таза пробиркаға 2 мл
натрий гидроксиді ерітіндісі құйылды және көк түсті мыс (II) гидроксиді
тұнбасы пайда болғанға дейін мыс(II) сульфатының ерітіндісі тамшылап
қосылды. Алынған қоспа шайқалып, оған фильтраттың бірінші жартысын
қосқанда, ашық көк түстің пайда болуы фильтратта лактоза бар екенін
көрсетті. Қыздырылған кезде ерітіндінің түсі қызылға өзгерді. Пробирка
қабырғаларында карамель иісі бар қоңыр тұнбаның пайда болуы лактозаның
болуын көрсетеді.

6. Сүт және сүт өнімдерінің құрамындағы С дәруменінің мөлшерін анықтау.
Аскорбин қышқылын анықтаудың йодометриялық әдісі қолданылды. Крахмал
ерітіндісі дайындалды: 1 г крахмалды суық суға сұйылтып, бір стакан
қайнаған суға құйып, шамамен 1 минут қайнатады. Йод ерітіндісі
дайындалды: йод концентрациясы 5\% болатын дәріханалық йод ерітіндісін
сумен 40 рет сұйылтып 0,125\% ерітінді алынды.20 мл сүт өлшеніп, 100 мл
көлемге дейін сумен сұйылтылды.1 мл крахмал ерітіндісін құйып, оған 1
тамшы йод ерітіндісі қосылды.10-15 с ішінде тұрақты көк түстің пайда
болуы сүтте аскорбин қышқылының болуын көрсетеді.

{\bfseries Нәтижелер мен талқылау.} Төмен лактозалы сүт тек жоғары сапалы
санаттағы сүттен жасалады. Алдымен сиыр сүтін тазартады, майлылығы 2,5\%
дейін май бойынша қалыпқа келтіреді, 15±2 МПа қысымда бір сатылы
гомогенизациядан өтеді, 76±2 °С температуралық режиммен 15-20 с
пастерленеді, одан кейін 4±2°С дейін салқындатылады.
\end{multicols}
\vspace{1em}
\begin{figs}[2-cурет. Төмен лактозалы сүт сусынын өндірудің технологиялық процесі]
    \fig[0.24\textwidth][\textwidth]{p/image30}[а]
    \fig[0.24\textwidth][\textwidth]{p/image31}[ә]
    \fig[0.24\textwidth][\textwidth]{p/image32}[б]
    \fig[0.24\textwidth][\textwidth]{p/image33}[в]
\end{figs}

\begin{multicols}{2}
Сүттің сапасын тексеріліп, кейін шикізат салқындатқыш арқылы үлкен
қалыпқа келтіру танктеріне түседі. Оған лактаза ферменті және үккіштен
өткізілген нори теңіз балдыры қосылады. Функционалдық ингредиент ретінде
нори теңіз балдырлары таңдалды, себебі олардың құрамында биологиялық
белсенді заттар - полисахаридтер, минералды элементтер және
антиоксиданттар бар. Теңіз балдырлары ас қорыту жүйесінің микрофлорасын
қалыпқа келтіруге ықпал ететін пребиотикалық қасиетке ие.

Төмен лактозалы сүтті сусын құрамына теңіз балдырларын қосу өнімнің
функционалдық бағытын күшейтумен қатар, оңтайлы мөлшерде қолданылған
жағдайда органолептикалық қасиеттерінің айтарлықтай нашарлауына
әкелмейді.

1 а, ә, б, в - суретте төмен лактозалы сүт сусынын өндірудің
технологиялық процесі көрсетілген.

Функционалдық қасиетті жоғарылату мақсатында сүтке теңіз балдырларының
тұнбалары қосылып дайындалғандығы туралы көптеген ғалымдардың
мақалаларынан көруге болады {[}7, 9{]}. Төмен лактозалы сүт сусынын
өндірудің қажетті мөлшері 1-кестеде көрсетілген.
\end{multicols}

\tcap{1-кесте. Өнімнің рецептурасы (л/100 г)}
\begin{longtblr}[
  label = none,
  entry = none,
]{
  cells = {c},
  cells = {font = \small},
  row{1} = {font=\bfseries\small},
  hlines,
  vlines,
}
Атаулары                  & Мөлшері \\
Сиыр сүті, 3,2 \%         & 1 л     \\
«Лактаза» ферменті        & 0,3 г   \\
Ұнтақталған теңіз балдыры & 0,2 г   
\end{longtblr}
\vspace{-1em}
\fig[0.5\textwidth]{p/image34}[3-сурет. Төмен лактозалы сүт сусынының 4 үлгісі]
\vspace{-1em}
\tcap{2-кесте. Төмен лактозалы сүт сусынының 4 үлгісін дайындау мөлшері}
\begin{longtblr}[
  label = none,
  entry = none,
]{
  cells = {c},
  cells = {font = \small},
  row{1} = {font=\bfseries\small},
  hlines,
  vlines,
}
Үлгі & Атаулары                  & Мөлшері, \% \\
№1   & Бақылау үлгісі            & -           \\
№2   & Ұнтақталған теңіз балдыры & 0,01        \\
№3   & Ұнтақталған теңіз балдыры & 0,02        \\
№4   & Ұнтақталған теңіз балдыры & 0,04        
\end{longtblr}

\begin{multicols}{2}
Лотте Дж. Кнудсен және т.б. ғалымдар сүтке лактаза ферментін қосып,
лактозаны гидролиздеу әдістерін сипаттаған {[}10{]}. Талдау үшін төмен
лактозалы сүт сусынының 4 үлгісі дайындалды. Бақылау үлгісі 100 мл сүтке
әртүрлі арақатынаста қосылды. Бақылау үлгісі ретінде фермент пен теңіз
балдыры қосылмаған пастерленген сиыр сүті (майлылығы 3,2 \%)
пайдаланылды. Бұл үлгі дайын өнімнің табиғи сипаттамаларымен салыстыру
үшін эталон ретінде алынды. Дайындалған үлгілердің мөлшері 2-кестеде
көрсетілген.

Барлық талдаулар кемінде үш қайталаумен жүргізілді. Әрбір үлгі
бойынша үш параллель анықтама алынып, орташа мән ± стандартты
ауытқу (SD) есептелді. Статистикалық өңдеу Excel бағдарламасының
Analysis ToolPak модулімен жүзеге асырылды. Айырмашылықтардың
сенімділігі Student t-критерийімен (p<0,05) анықталды.
Зертханалық жағдайда тәжірибе үлгісі мен бақылау үлгілері алынып,
талданды. Өнімнің органолептикалық және физика-химиялық көрсеткіштері
сыртқы түрі, консистенция, иіс пен дәм бойынша салыстыра анықталды.3-ші
кестеде төмен лактозалы сүт сусынының органолептикалық бағалау
нәтижелері келтірілген.
\end{multicols}

\tcap{3-кесте. Төмен лактозалы сүт сусынының органолептикалық бағалау нәтижелері}
\begin{longtblr}[
  label = none,
  entry = none,
]{
  width = \linewidth,
  colspec = {Q[169]Q[137]Q[210]Q[210]Q[213]},
  cells = {c},
  cells = {font = \small},
  cell{1}{1} = {font=\bfseries\small},
  cell{1}{2} = {font=\bfseries\small},
  hlines,
  vlines,
}
Көрсеткіш атауы                 & Бақылау үлгісі           & № \textbf{1}                             & № \textbf{2}                             & № \textbf{3}                             \\
Консистенциясы және сыртқы түрі &                          & Біртекті теңіз балдырының бөлшектері бар & Біртекті теңіз балдырының бөлшектері бар & Біртекті теңіз балдырының бөлшектері бар \\
Дәмі және иісі                  & Төмен лактозалы сүт дәмі & Балдырдың дәмі сезілмеді                 & Балдырдың дәмі кішкене сезіледі          & Балдырдың дәмі қатты,дәмі ауызда қалады  \\
Түсі                            & Ақ түсті                 & Ақ түсті                                 & Ақ түсті                                 & Ақ түсті,балдыр ұнтақтары көрінеді       
\end{longtblr}

\begin{multicols}{2}
Дайын өнімнің органолептикалық көрсеткіштері С. Сейфуллин атындағы Қазақ
агротехникалық университетінің Тамақ және қайта өңдеу өндірісінің
технологиясы кафедрасының меңгерушісіне, оқытушыларына дегустация
жүргізіп, 0,02 \% қатынасында дайындалған төмен лактозалы сүт сусынын
оңтайлы деп анықталды. Себебі зерттеу нәтижелері бойынша теңіз
балдырларын 0,02 \% мөлшерде қосу ең оңтайлы нұсқа екені анықталды. Бұл
дозировкада сусын жоғары органолептикалық көрсеткіштермен, тұрақты
физика-химиялық параметрлермен және жақсы құрылымдық қасиеттермен
сипатталды. Теңіз балдырларының 0,01 \% мөлшері функционалдық әсердің
жеткіліксіздігімен ерекшеленсе, 0,04 \% деңгейінде қосу сусынның дәмі
мен консистенциясының нашарлауына және қышқылдықтың шамадан тыс артуына
алып келді. Айырмашылықтардың сенімділігі --~\emph{p <{}
0,05}~деңгейінде есептелді.

Зерттелетін сусын үлгілері (4 ± 2) °C температурада 7 тәулік бойы
сақталды. Зерттеулер сақтау мерзімінің 1-, 3- және 7-күндері жүргізілді.
Органолептикалық бағалау 7 сарапшыдан тұратын дегустациялық комиссиямен
жүргізілді. Бағалау бейтарап жарықтандырылған бөлмеде, ауа температурасы
20-22 °C жағдайында өткізілді. Сусынның дәмі, иісі, түсі және
консистенциясы 5 балдық шкала бойынша бағаланды. Барлық тәжірибелер үш
реттік қайталанымда орындалды. Алынған нәтижелер вариациялық статистика
әдістері арқылы өңделіп, орташа мәні мен стандартты ауытқуы есептелді.

Зерттеу барысында органолептикалық, физика-химиялық және
микробиологиялық көрсеткіштерді таңдау функционалдық сүтті сусындарға
қойылатын талаптарға негізделді. Органолептикалық көрсеткіштер (дәмі,
иісі, түсі және консистенциясы) өнімнің тұтынушылық қасиеттерін
бағалауда негізгі рөл атқарады. Қышқылдықты (°Т) анықтау ашыту
процесінің қарқындылығын, өнімнің тұрақтылығын және сақтау барысында
сапасының өзгеруін бағалауға мүмкіндік береді. Лактоза мөлшерін анықтау
ашыту тиімділігін растау және өнімнің «лактозасы төмен» санатына
сәйкестігін дәлелдеу үшін жүргізілді.

Зерттелген көрсеткіштер функционалдық сүт өнімдерінің қауіпсіздігі мен
сапасына қойылатын талаптарды регламенттейтін ТР ТС 033/2013 «Сүт және
сүт өнімдерінің қауіпсіздігі туралы» техникалық регламентіне сәйкес
таңдалды. Кестелерде көрсетілген деректер өнім үлгілерінің барлық мәндер
үшін МЕМСТ Р 53590-2009 және МЕМСТ 31761-2012 сәйкестігін көрсетеді.4 -
кестеде теңіз балдыры қосылған төмен лактозалы сүттің органолептикалық
көрсеткіштері көрсетілген.
\end{multicols}

\tcap{4-кесте. Теңіз балдыры қосылған төмен лактозалы сүттің
органолептикалық көрсеткіштері}
\begin{longtblr}[
  label = none,
  entry = none,
]{
  cells = {c},
  cells = {font = \small},
  row{1} = {font=\bfseries\small},
  hlines,
  vlines,
}
Көрсеткіштер                   & №2 тәжірибелік үлгісінің жалпы сипаттамасы \\
Сыртқы түрі мен консистенциясы & Біркелкі                                   \\
Дәмі мен иісі                  & Таза, сүтке тән дәм                        \\
Түсі                           & Ақ түсті, біркелкі                         
\end{longtblr}

Теңіз балдыры қосылған төмен лактозалы сүттің органолептикалық
көрсеткіштеріне тоқталсақ, өнімнің сыртқы түрі біркелкі, консистенциясы
біртекті, иіс пен дәм бойынша сүтке тән дәм, ащы емес, біркекі ақ түске
ие болды.5 - кестеде теңіз балдыры қосылған төмен лактозалы сүттің
физика-химиялық көрсеткіштері көрсетілген.

\tcap{5-кесте. Теңіз балдыры қосылған төмен лактозалы дайын сүттің
физика-химиялық көрсеткіштері}
\begin{longtblr}[
  label = none,
  entry = none,
]{
  width = \linewidth,
  colspec = {Q[150]Q[350]Q[200]},
  cells = {c},
  cells = {font = \small},
  row{1} = {font=\bfseries\small},
  hlines,
  vlines,
}
Көрсеткіштер                    & МЕМСТ 31450-2013 стандарты бойынша майлылығы 1,5 \% сиыр сүті & {\textbf{Теңіз балдыры қосылған}\\\textbf{төмен лактозалы сүті}} \\
Тығыздығы,/м\textsuperscript{3} & 1027-1029 /м\textsuperscript{3}                               & 28,77 /м\textsuperscript{3} ± 0,12                               \\
Құрғақ зат, \%                  & 8,2 \%                                                        & 8,37 \%                                                          \\
Қышқылдығы, °Т                  & 21 °Т                                                         & 19 °Т ± 0,4                                                      \\
Майлылығы, г                    & 1,5 г                                                         & 1,1 г ± 0,1                                                      \\
Қосылған су, г/л                & 1,5 г/л                                                       & 0 г/л ± 0,01                                                     \\
Ақуыздар, г                     & 3,0 г                                                         & 3,09 г ± 0,001                                                   \\
Температура, °С                 & 4 ± 2 °С                                                      & 4 ± 2 °С ± 0,01                                                  \\
Лактоза, г                      & 4,9 г                                                         & 0,1 г ± 0,1                                                      
\end{longtblr}

\begin{multicols}{2}
Кестедегі зерттеу нәтижесі бойынша зерттеу үлгісінде МЕМСТ 31450-2013
стандарты бойынша майлылығы 1,5 \% сиыр сүті мен теңіз балдыры қосылған
төмен лактозалы сүттің физика-химиялық көрсеткіштері көрсетілген.
Ферментация процесі барысында сусындағы лактоза мөлшерінің айтарлықтай
төмендегені байқалды, бұл қолданылған ферменттің тиімділігін дәлелдейді.
Ашыту аяқталғаннан кейін өнімнің лактоза мөлшері төмен лактозалы сүт
өнімдеріне қойылатын талаптарға сәйкес келді.

Теңіз балдырлары мөлшерінің артуы титрленетін қышқылдықтың біртіндеп
жоғарылауына әсер етті, бұл олардың биологиялық белсенді
компоненттерінің ашыту процесіне ықпал етуімен түсіндіріледі.
Нәтижесінде теңіз балдыры қосылған төмен лактозалы сүттің құрамындағы
лактоза мөлгері 0,1 г тең, бұл лактозаға төзбеушілікпен ауыратын адамдар
үшін ұсынылған нормаға сәйкес келуін дәлелдейді.

Төмен лактозалы сүт сусынын өндіру технологиялық схемасы келесі негізгі
сатылардан тұрады:
\end{multicols}

\fig[0.5\textwidth]{p/image48}[3-сурет. Төмен лактозалы сүт сусынының
4 үлгісі]

\begin{multicols}{2}
{\bfseries Қорытынды.} Зерттеу тақырыбы бойынша ғылыми-техникалық
әдебиеттер мен патенттік ақпараттар талданды және жүйеленді.
Функционалдық бағытта шикізат таңдалды. Құрамы мен қасиеттері зерттелді.
Төмен лактозалы сүт сусынының рецептурасы әзірленді. Зерттеу нәтижелері
төмен лактозалы сүт сусынын өндіруде 0,02 \% теңіз балдыры мен лактаза
ферментін пайдалану өнімнің органолептикалық және физика-химиялық
қасиеттерін жақсартатынын көрсетті. Алынған нәтижелер гипотезаны
растады: теңіз балдырлары өнімнің құрылымдық және тағамдық
сипаттамаларын тұрақтандырды, ал лактаза ферменті лактозаны 98 \%-ға
дейін ыдыратты. Болашақта өнімнің сақтау мерзімін, микробиологиялық
тұрақтылығын және тұтынушылардың қабылдау деңгейін зерттеу жоспарлануда.
\end{multicols}

\begin{center}
{\bfseries Әдебиеттер}
\end{center}

\begin{refs}
1. Дмитриева И.Е., Иванова Л.А., Донская Г.А. Разработка синбиотического
кисломолочного продукта с биологически активной растительной добавкой.
Биотехнология и продукты биоорганического синтеза: Сборник материалов
национальной научно-практической конференции/ ФГБОУ ВО «Московский
государственный университет пищевых производств». - 2018. -С.189-192.

2. Антипова Т.А., Фелик С.В., Симоненко С.В., Кудряшова О.В. Получение
низколактозного молока для специализированных продуктов детского питания
/ Т. А. Антипова, // Пищевая промышленность. -2020. -№ 10.- С.41-46.
DOI 10.24411/0235-2486-2020-10105.

3. Данильчук Т.Н., Ганина В.И., Головин М.А. Низколактозные молочные
продукты пути получения / Молочная промышленность. - 2013. -№ 11. - С.
41-42.

4. Трушина Ю.В. Применение фермента «Lacta-free» при получении
низколактозных продуктов / Молочнакя промышленность. -2014. -№ 11. -С.
31-32.

5. Knudsen L.J., Rauh V., Pedersen J.N., Dekker P., Otzen D.E., Larsen
L.B., Nielsen S.D. Insight into protein crosslinking and casein
polymerization in pre- and posthydrolyzed lactose-free
ultra-high-temperature milk during long-term storage//Journal of Dairy
Science. -2024. -Vol.107(12). -P.10497-10511.
\href{https://doi.org/10.3168/jds.2024-25103}{DOI
10.3168/jds.2024-25103}.

6. Titov E.I., Tikhomirova N.A., Nguyen B.C., et al.~Research of lactose
hydrolysis depending on the type of the enzyme~//~IOP Conference Series:
Earth and Environmental Science. -2020. -Vol.548 (8):082040.
\href{https://doi.org/0.1088/1755-1315/548/8/082040}{DOI
10.1088/1755-1315/548/8/082040}.

7. O' Sullivan A. M. et al. An examination of the
potential of seaweed extracts as functional ingredients in milk//
International Journal of Dairy Technology. -2014. --Vol.67(2). -P.
182-193. \href{https://doi.org/10.1111/1471-0307.12121}{DOI
10.1111/1471-0307.12121}.

8. Orzuna-Orzuna J.F., Lara-Bueno A., Mendoza-Martínez G. D.,
Miranda-Romero L. A., Vázquez Silva G., de la Torre-Hernández M. E.,
Sánchez-López N., Hernández-García P. A.~Meta-Analysis of Dietary
Supplementation with Seaweed in Dairy Cows: Milk Yield and Composition,
Nutrient Digestibility, Rumen Fermentation, and Enteric Methane
Emissions~//~Dairy. -2024. -Vol.5(3). -P.464-479.
\href{https://doi.org/10.3390/dairy5030036}{DOI 10.3390/dairy5030036}.

9. Mummah S., Oelrich B., Hope J., Vu Q., Gardner C.D.~Effect of raw
milk on lactose intolerance: a randomized controlled pilot
study~//~Annals of Family Medicine. -2014. -Vol.12(2). -P.134-141.
\href{https://doi.org/10.1370/afm.1618}{DOI 10.1370/afm.1618}.

10. Knudsen L.J., Nielsen S.D-H., Dekker P., Otzen D.E., Rauh V., Larsen
L.B. Impact of hydrolysis method and lactase preparation on proteolysis
and glycation in long-term stored lactose-hydrolysed UHT milk //
International Dairy Journal. -2024. -Vol.154:105946.
\href{https://doi.org/10.1016/j.idairyj.2024.105946}{DOI
10.1016/j.idairyj.2024.105946}.
\end{refs}

\begin{center}
{\bfseries References}
\end{center}

\begin{refs}
1. Dmitrieva I.E., Ivanova L.A., Donskaja G.A. Razrabotka
sinbioticheskogo kislomolochnogo produkta s biologicheski aktivnoj
rastitel' noj dobavkoj. Biotehnologija i produkty
bioorganicheskogo sinteza: Sbornik materialov
nacional' noj nauchno-prakticheskoj konferencii/ FGBOU VO
«Moskovskij gosudarstvennyj universitet pishhevyh proizvodstv». - 2018.
-S.189-192. {[}in Russian{]}

2. Antipova T.A., Felik S.V., Simonenko S.V., Kudrjashova O.V.
Poluchenie nizkolaktoznogo moloka dlja specializirovannyh produktov
detskogo pitanija / T. A. Antipova, // Pishhevaja
promyshlennost'. -2020. -№ 10.- S.41-46. DOI
10.24411/0235-2486-2020-10105. {[}in Russian{]}

3. Danil' chuk T.N., Ganina V.I., Golovin M.A.
Nizkolaktoznye molochnye produkty puti poluchenija / Molochnaja
promyshlennost'. - 2013. -№ 11. - S.41-42. {[}in
Russian{]}

4. Trushina Ju.V. Primenenie fermenta «Lacta-free» pri poluchenii
nizkolaktoznyh produktov / Molochnakja promyshlennost'.
-2014. -№ 11. -S.31-32. {[}in Russian{]}

5. Knudsen L.J., Rauh V., Pedersen J.N., Dekker P., Otzen D.E., Larsen
L.B., Nielsen S.D. Insight into protein crosslinking and casein
polymerization in pre- and posthydrolyzed lactose-free
ultra-high-temperature milk during long-term storage//Journal of Dairy
Science. -2024. -Vol.107(12). -P.10497-10511.
\href{https://doi.org/10.3168/jds.2024-25103}{DOI
10.3168/jds.2024-25103}.

6. Titov E.I., Tikhomirova N.A., Nguyen B.C., et al.~Research of lactose
hydrolysis depending on the type of the enzyme~//~IOP Conference Series:
Earth and Environmental Science. -2020. -Vol.548 (8):082040.
\href{https://doi.org/0.1088/1755-1315/548/8/082040}{DOI
10.1088/1755-1315/548/8/082040}.

7. O' Sullivan A. M. et al. An examination of the
potential of seaweed extracts as functional ingredients in milk//
International Journal of Dairy Technology. -2014. --Vol.67(2). -P.
182-193. \href{https://doi.org/10.1111/1471-0307.12121}{DOI
10.1111/1471-0307.12121}.

8. Orzuna-Orzuna J.F., Lara-Bueno A., Mendoza-Martínez G. D.,
Miranda-Romero L. A., Vázquez Silva G., de la Torre-Hernández M. E.,
Sánchez-López N., Hernández-García P. A.~Meta-Analysis of Dietary
Supplementation with Seaweed in Dairy Cows: Milk Yield and Composition,
Nutrient Digestibility, Rumen Fermentation, and Enteric Methane
Emissions~//~Dairy. -2024. -Vol.5(3). -P.464-479.
\href{https://doi.org/10.3390/dairy5030036}{DOI 10.3390/dairy5030036}.

9. Mummah S., Oelrich B., Hope J., Vu Q., Gardner C.D.~Effect of raw
milk on lactose intolerance: a randomized controlled pilot
study~//~Annals of Family Medicine. -2014. -Vol.12(2). -P.134-141.
\href{https://doi.org/10.1370/afm.1618}{DOI 10.1370/afm.1618}.

\href{https://doi.org/10.3390/dairy5030036\%209}{10}. Knudsen L.J.,
Nielsen S.D-H., Dekker P., Otzen D.E., Rauh V., Larsen L.B. Impact of
hydrolysis method and lactase preparation on proteolysis and glycation
in long-term stored lactose-hydrolysed UHT milk // International Dairy
Journal. -2024. -Vol.154:105946.
\href{https://doi.org/10.1016/j.idairyj.2024.105946}{DOI
10.1016/j.idairyj.2024.105946}.
\end{refs}

\begin{info}
\hspace{1em}\emph{{\bfseries Авторлар туралы мәлімет}}

Жасқайрат Ш.Ж. - т.ғ.м., докторант, С.Сейфуллин атындағы Қазақ
агротехникалық зерттеу университеті, Астана, Қазақстан, e-mail:
shynarai\_92@mail.ru;

Машанова Н.С. - техника ғылымдарының докторы, аға оқытушы, С.Сейфуллин
атындағы Қазақ агротехникалық зерттеу университеті, Астана, Қазақстан,
e-mail:
nurmashanova@gmail.com;

Жаксыгалиева Д. С. - т.ғ.м., аға оқытушы, Жәңгір хан атындағы Батыс
Қазақстан аграрлық технологиялық университеті, Орал, Қазақстан, e-mail:
dariga\_aidos@mail.ru;

Бекболатова М.Е. - т.ғ.м., докторант, С.Сейфуллин атындағы Қазақ
агротехникалық зерттеу университеті, Астана, Қазақстан, e-mail:
mariambekbolatova93@gmail.com.

\hspace{1em}\emph{{\bfseries Information about the authors}}

Zhaskairat Sh. Zh. - PhD student, Saken Seifullin Kazakh Agrotechnical
Research University, Astana, Kazakhstan, e-mail:
shynarai\_92@mail.ru;

Mashanova N.S. - Doctor of Technical Sciences, Senior Lecturer, Saken
Seifullin Kazakh Agrotechnical Research University, Astana, Kazakhstan,
e-mail: nurmashanova@gmail.com;

Zhaksygalieva D.S. - Senior Lecturer, Master of Technical Sciences,
Zhangir Khan West-Kazakhstan agrarian technical University, Uralsk,
Kazakhstan, e-mail:
dariga\_aidos@mail.ru;

Bekbolatova M.E. - PhD student, Saken Seifullin Kazakh Agrotechnical
Research University, Astana, Kazakhstan, e-mail:
mariambekbolatova93@gmail.com.
\end{info}
